\section{Introduction}

Error-correcting codes (ECCs) are fundamental to protecting memory systems from data corruption. Memory ECCs are commonly differentiated by correction and detection capability, redundancy, and decoding complexity. Those properties establish what a code can protect, but implementation decisions determine whether that protection is practical within area, clock, latency, routing, and energy constraints.

Algorithm-level descriptions do not expose standard-cell area, achievable timing after routing, interconnect burden, or activity-dependent energy. Moreover, the mapping from a code to hardware is not unique: two architectures can implement exactly the same coding semantics while inserting different register boundaries and logic structures. Hardware-aware SECDED construction and checker optimization are well established~\cite{hsiao1970,ghosh2004,ghosh2005,reviriego2013}, as are cross-layer studies that trade ECC strength against memory cost~\cite{delbel2014,pajouhi2015}. This makes a controlled implementation comparison more useful than a claim that physical cost has been overlooked.

We use the workflow in Fig.~\ref{fig:method}: correctness and implementation identity are qualified before physical comparison; guaranteed reliability semantics remain separate from observations outside that guarantee; accepted designs are implemented under common constraints; and only then are reliability and physical cost interpreted together. The procedure is fail-closed: unsupported or incomplete measurements remain unavailable rather than being imputed.

The strongest controlled contrast is between combinational and pipelined extended-Hamming SECDED (72,64) RTL. Exact equivalence after temporal alignment fixes their coding semantics. Pipelining increases area by \SecAreaPct\% and detailed wirelength by \SecWirePct\%, and adds two cycles of request latency, while increasing achieved Fmax by \SecFmaxPct\% and reducing equal-work, steady-stream energy per operation by \SecEnergyReductionPct\%. A second contrast is the evaluated BCH (78,64,$t=2$) RTL: relative to combinational SECDED, it uses \BchAreaPct\% more area and \BchWirePct\% more wirelength, achieves \BchFmaxReductionPct\% lower Fmax, and fails the common 10 ns target with \BchWns\,ns worst negative slack (WNS).

This paper makes three contributions:

\begin{enumerate}
    \item an identity-preserving cross-layer evaluation method that admits correctness-qualified ECC implementations and keeps guarantees, observations, architecture, and unavailable physical data distinct;
    \item a quantified non-dominated post-route trade-off between two exact-equivalent SECDED microarchitectures; and
    \item an implementation-scoped correction-versus-feasibility result for the evaluated BCH RTL, while retaining unavailable Hsiao PPA explicitly rather than estimating it.
\end{enumerate}

All findings are bounded to the evaluated RTL, SKY130HD environment, and frozen experimental conditions.

